% ============================================================================
% SEC01.TEX - Section 4.1: Setup Project Space Survivor
% ============================================================================

\section{Setup Project Space Survivor}

\begin{learningoutcomes}
\item Membuat Project 2D baru di Unity Hub
\item Mengatur struktur folder project yang rapi
\item Mengimpor asset dasar untuk game Space Survivor
\end{learningoutcomes}

\subsection{Membuat Project Baru}

Langkah pertama dalam membuat game \textit{Space Survivor} adalah membuat project baru:

\begin{enumerate}
    \item Buka Unity Hub.
    \item Klik tombol \textbf{New Project}.
    \item Pilih template \textbf{2D (URP)} atau \textbf{2D Core}. Karena game ini grafisnya sederhana, 2D Core sudah cukup.
    \item Beri nama project: \texttt{SpaceSurvivor\_[NIM\_ANDA]}.
    \item Klik \textbf{Create Project} dan tunggu hingga Unity Editor terbuka.
\end{enumerate}

\subsection{Mengatur Folder Project}

Agar project tetap rapi, buatlah folder-folder berikut di dalam panel \textbf{Project} (Assets):

\begin{itemize}
    \item \texttt{Sprites}: Untuk menyimpan gambar pesawat, asteroid, dan background.
    \item \texttt{Scripts}: Untuk menyimpan kode C\# (PlayerController, dll).
    \item \texttt{Prefabs}: Untuk menyimpan objek yang akan digunakan berulang kali (Bullet, Asteroid).
    \item \texttt{Scenes}: Untuk menyimpan level game.
    \item \texttt{Audio}: Untuk suara tembakan dan ledakan.
\end{itemize}

\subsection{Impor Aset}

Untuk latihan ini, kita akan menggunakan \textit{placeholder assets} (bentuk geometri sederhana) yang bisa dibuat langsung di Unity, atau Anda bisa mengunduh aset "Space Shooter" gratis dari \textit{Kenney.nl}.

\begin{tips}
Untuk sementara, kita akan menggunakan \textbf{Sprite Shape} bawaan Unity (Square, Circle, Triangle) sebagai representasi pesawat dan asteroid.
\end{tips}
